\thispagestyle{plain}
\begin{center}
{\fontsize{10}{12}\selectfont Manuscrito de recopilación para transferencia académica\par}
\vspace{1.5mm}
{\fontsize{17}{20.5}\selectfont\bfseries Relaciones estructurales\par entre las constantes físicas\par}
\vspace{1.5mm}
{\fontsize{11}{13.5}\selectfont Acción, reciprocidad geométrica\par y memoria de las transformaciones\par}
\vspace{1.5mm}
{\fontsize{11.5}{14}\selectfont Oumar Haidara Fall \textperiodcentered\ Rubén Ramos Balsa\par}
\vspace{1.5mm}
{\fontsize{8}{10}\selectfont Coautoría sin prelación de contribución\par}
\vspace{1.5mm}
{\fontsize{8}{10}\selectfont Integración y recopilación documental generadas por ChatGPT - Astra.\par}
\end{center}
\vspace{1.5mm}
\begingroup\fontsize{9.5}{10.7}\selectfont\setlength{\parskip}{1.5pt}\renewenvironment{abstract}{\begin{center}\bfseries\abstractname\end{center}\noindent}{\par}\begin{abstract}
Se establecen relaciones de compatibilidad, reciprocidad y reconstrucción
entre las estructuras de acción, masa y respuesta constitutiva obtenidas
mediante Holografía Modular Triádica y Mecánica Dimensional. El desarrollo
parte de la composición de la estructura aritmética de caminos
\emph{All Possible Paths} (APP), la firma ternaria orientada TRIT y el
\emph{Topological Phase Kernel} (TPK), que actualiza y transporta el estado
enriquecido. Esta composición conserva orientación, acarreo, memoria y
retorno. Las constantes intervienen como evaluaciones de las
estructuras previamente construidas; su composición permite determinar qué
propiedades comparten y qué información recupera una observación conjunta.
En la colección constitutiva normalizada se demuestra que la velocidad de
propagación y el funcional de Barbero--Immirzi determinan los autovalores
etiquetados de los dos canales. Los logaritmos de velocidad e impedancia
constituyen el gradiente de un potencial espectral estrictamente cóncavo.
Su Hessiano caracteriza la sensibilidad de la inversión y proporciona cotas
explícitas de estabilidad. El potencial comparte un momento espectral de
profundidad dos con la construcción orientada cuyo valor extremo es Catalán;
la respuesta constitutiva restituye esa colección de manera única con su
condición de frontera en el origen.

La composición de cuartos de giro asociados a formas elípticas conjugadas
produce un transporte relativo unimodular cuyas dilataciones recíprocas
permiten reconstruir el coeficiente adimensional de acción, conservadas la
coordenada de clausura y la orientación. El levantamiento dodecafásico y la
iteración del transporte caracterizan la memoria del orden de las
operaciones. El análisis cronológico distingue el crecimiento de la
información de la tasa entrópica por posición, y la reentrada de memoria
determina su contribución a la respuesta operatoria. El retorno inscrito y
su composición polar determinan una longitud $L_\alpha$. Conservar
conjuntamente acción y área fija $G=c^3L_\alpha^2/\hbar$ en la realización
radial, también bajo variaciones no conmutativas y reducción compatible
de memoria.

La memoria electrónica se restituye mediante invariantes lineales y
espectrales que conservan el orden del registro; el transporte entre
multisecciones admite un criterio explícito de conservación de masa.
La composición térmica reúne capacidades, multiplicidades y memoria
espectral en una traza exacta, y construye su respuesta termométrica
para una lectura marcada con referencia fija. Esta realización enlaza
energía, información, acción y área, conservando las condiciones de
normalización y de observación de cada magnitud.
La composición transportada de las respuestas y la reconstrucción
Higgs--impedancia enlazan los canales constitutivos con el sector
electrodébil. La sustitución electrónica conserva su escala en Fermi,
la cancela en los cocientes angulares y transporta la longitud inversa
mediante una congruencia operatoria, también sin conmutación.

Finalmente, bajo las
condiciones circular y térmica especificadas, una misma coordenada de masa
determina la frecuencia propia, dos longitudes recíprocas de producto
invariante y la ponderación térmica del ciclo. Los resultados reúnen así
producción, compatibilidad y restitución en una descripción matemática
común, con demostraciones explícitas y con la distinción entre cambio de
representación, transporte relativo y realización física.


El lector dodecafásico uniforme determina además una sucesión de primeras raíces de retorno crítico. Se demuestra que su cociente de separaciones converge a la constante de Feigenbaum, conservando la selección original, la orientación y la memoria de la construcción.

% Adición al resumen de VII; no sustituye su contenido anterior.
La misma composición se prolonga a una acción total de geometría,
calibre y materia: incorpora conjuntamente la gravitación y las
interacciones fuerte, débil y electromagnética, conservando las
representaciones quirales y los acoplamientos previamente generados.
Se obtienen su corriente de espín, la fuente métrica total y los
términos torsionales cruzados entre especies. La transformación de
Legendre y las identidades de simetría producen restricciones
conjuntas de primera clase y su propagación. Sobre su superficie
regular se calcula la curvatura de la conexión canónica total y se
demuestra su anulación en las direcciones de deformación
características. El resultado conserva las condiciones de borde y
la posible holonomía global del estado enriquecido.



\par\smallskip\noindent\textbf{Palabras clave:} acción; reciprocidad geométrica; memoria operatoria; constantes físicas; Holografía Modular Triádica; Mecánica Dimensional.
\end{abstract}


\endgroup
